% Generated draft: requires mathematical and source review.
\documentclass[11pt]{article}
\usepackage[margin=25mm]{geometry}
\usepackage{fontspec}
\setmainfont{Latin Modern Roman}
\usepackage{amsmath,amssymb,mathtools,unicode-math}
\setmathfont{Latin Modern Math}
\usepackage{enumitem,needspace,xurl,hyperref}
\setlength{\parindent}{0pt}
\setlength{\parskip}{5pt}
\setlength{\emergencystretch}{3em}
\newcommand{\N}{\mathbb{N}}
\newcommand{\Z}{\mathbb{Z}}
\newcommand{\Q}{\mathbb{Q}}
\newcommand{\R}{\mathbb{R}}
\newcommand{\C}{\mathbb{C}}
\newcommand{\F}{\mathbb{F}}
\newcommand{\A}{\mathbb{A}}
\newcommand{\PP}{\mathbb{P}}
\newcommand{\E}{\mathbb{E}}
\newcommand{\eps}{\varepsilon}
\newcommand{\dd}{\,\mathrm{d}}
\begin{document}
\section*{UnsolvedMath ID: 30000264}
\textbf{Source problem code:} catalogue entry
\subsection{Definitions and mathematical statement}
The source specifies the subset
\[
S=\{u\in\R^p\setminus\{0\}:(Au,u)\leq0\}
\]
and asks for its topology in the setting of sign-changing Yamabe asymptotes. The origin is excluded, the zero level is included, and no positivity restriction on the coordinates of \(u\) is stated. The set carries the subspace topology from \(\R^p\).

If \(A\) denotes a fixed real linear operator on \(\R^p\) and the parentheses denote the Euclidean inner product, the literal object is
\[
S_A=\{u\ne0:u^{\mathsf T}Au\leq0\},
\qquad
u^{\mathsf T}Au=u^{\mathsf T}\frac{A+A^{\mathsf T}}2u.
\]
In this interpretation only the symmetric part contributes, and the set is invariant under positive scaling. Its radial description is
\[
S_A\cong(0,\infty)\times\{v\in S^{p-1}:v^{\mathsf T}Av\leq0\},
\qquad u\longmapsto(\|u\|,u/\|u\|).
\]
These observations specify the literal quadratic-set reading; they do not recover which operator, parameter dependence, or extra asymptotic information the source intends. In particular, a question about one fixed quadratic cone must not be identified without explanation with a possibly richer family of Yamabe asymptotic sets.

The supplied record does not define \(A\), \(p\), the inner product, or the relation to the Yamabe problem. It also does not say whether ``determine the topology'' requests a homeomorphism type, homotopy type, homology, connectedness, or behavior in a parameter family. The precise source data and target invariant are therefore unresolved. The supplied assessment checked 2026-08-21 reports that the report calls the topology problem open and that no resolution in that formulation was located; this dated label is not independently certified for the underspecified quadratic-set interpretation.
\subsection{Short English statement}
Determine the topology of the punctured nonpositive set defined by \((Au,u)\), in the intended Yamabe asymptotic setting. The operator and its geometric role are missing; the literal fixed-matrix reading is kept separate from an unspecified richer asymptotic problem.
\subsection{Sources}
\begin{itemize}
\item Local UnsolvedMath registry, permanent ID 30000264, catalogue entry, \emph{Topology of Yamabe Asymptotic Sets}.
\item \emph{Topological and Variational Methods for Differential Equations} (2006), Oberwolfach Reports, pp. 1601--1678, DOI \url{https://doi.org/10.4171/owr/2005/29}.
\item Supplied report-text link identifying the topology problem: \url{https://oa.tib.eu/renate/bitstreams/91dd0822-361f-44ee-a77f-fa862bb21ab6/download}.
\item Supplied literature triage checked 2026-08-21; no definition of the operator or separate theorem citation is reproduced.
\end{itemize}
\end{document}
